\documentclass[12pt]{article}

\usepackage[utf8]{inputenc}
\usepackage[T2A]{fontenc}
\usepackage[english,russian]{babel}
\usepackage[margin=1in]{geometry}
\usepackage{amsmath}
\usepackage{amssymb}

\setlength{\parindent}{0pt}
\setlength{\parskip}{1\baselineskip}

\begin{document}

\section*{Нормализация лямбда-терма}

\text{Исходный терм:}
\[
  ((\lambda a. ((\lambda b. b\ b)\ (\lambda b. b\ b)))\ b)
  \ ((\lambda c. c\ b)\ (\lambda a. a)).
\]

Обозначим:
\[
  \Omega = (\lambda b. b\ b)\ (\lambda b. b\ b).
\]

Тогда исходный терм можно преобразовать следующим образом:
\[
\begin{aligned}
&
((\lambda a. ((\lambda b. b\ b)\ (\lambda b. b\ b)))\ b)
\ ((\lambda c. c\ b)\ (\lambda a. a))
\\[8pt]
&\to_{\beta}
((\lambda b. b\ b)\ (\lambda b. b\ b))
\ ((\lambda c. c\ b)\ (\lambda a. a))
\\[8pt]
&=
\Omega\ ((\lambda c. c\ b)\ (\lambda a. a))
\\[8pt]
&\to_{\beta}
\Omega\ ((\lambda a. a)\ b)
\\[8pt]
&\to_{\beta}
\Omega\ b.
\end{aligned}
\]

где
\[
  \Omega = (\lambda b. b\ b)\ (\lambda b. b\ b).
\]

\end{document}